\section{Thảo luận}\label{sec:discussion}

\subsection{Diễn giải các phát hiện chính}

\textbf{Khi đề đã chi tiết, quy trình chủ yếu thêm chi phí.} Với đề nêu đủ tên, chữ ký và quy tắc, làm theo đặc
tả không làm tác tử đạt test ẩn nhiều hơn: H1 chưa kết luận được, nhưng KTC không chứa giá trị dương nào, còn
chi phí tăng rõ ở mọi ô (H2). Với hai mô hình Claude, hiệu số âm hoặc bằng không ở gần như mọi tầng. Cách giải
thích hợp lý nhất là phần lợi của bản thiết kế, tức nêu rõ giao diện, trường hợp biên và chia bước
\cite{zi_2025}, đã nằm sẵn trong đề bài. Dấu hiệu có ích duy nhất nằm ở Codex trên ba nhiệm vụ tầng lớn; với ba
nhiệm vụ và H5 chưa kết luận được, đây là giả thuyết cho nghiên cứu sau, không phải bằng chứng rằng quy trình
``trả vốn'' khi việc đủ lớn.

\textbf{Phần lớn lượt hỏng đến từ thông tin đề không có.} \NumLabOutside{} trên \NumLabRuns{} lượt hỏng của
Claude không do logic tác tử tự viết sai, mà do test ẩn kiểm một chi tiết đề không chốt hoặc do nối sai với mã có
sẵn, và hai loại này chiếm tỉ lệ như nhau ở hai cách làm. Bước thiết kế không tự biết chi tiết mà đề bỏ ngỏ: như
ca \texttt{f15}, nó chốt một cách đoán, rồi test viết trước kiểm lại chính cách đoán đó. Thông tin còn thiếu chỉ
người giao việc có, nên quy trình tự động không bù được.

\textbf{Khi đề sơ sài, độ rõ của đề quan trọng hơn quy trình.} E2 gợi ý rằng bỏ tên và chữ ký làm điểm hành vi
giảm mạnh ở cả hai cách, lớn hơn nhiều mọi khác biệt giữa hai cách. Quy trình bù lại vài điểm ở việc nhỏ, rõ
nhất với Sonnet, có lẽ vì bước thiết kế buộc tác tử tự nêu điều đề không nói. Vì E2 thêm sau khi thấy kết quả
chính và tương tác có KTC chứa 0, đây là xu hướng gợi ý, chưa phải kết luận đã kiểm.

\textbf{Giá trị chưa được đo của quy trình là bộ test để lại.} Nhánh theo đặc tả để lại test ở phần lớn số lượt,
trong khi nhánh vibe của Sonnet và Codex gần như không. Test ẩn không đo giá trị này; nó chỉ hiện ra ở những lần
sửa sau, ngoài phạm vi thí nghiệm.

\subsection{Đối chiếu với tài liệu}

Kết quả khớp với nghiên cứu quan sát lớn nhất: Hill \cite{hill_2026} không thấy pull request có tài liệu đặc tả
đưa vào ít lỗi hơn; ở đây nhánh theo đặc tả cũng không đạt nhiều test ẩn hơn và không ít hồi quy hơn với hai mô
hình Claude. Phần hồi cứu (RQ2) cùng hướng: commit có hay không có plan bị sửa lại gần như ngang nhau. Đánh đổi
mà Agarwal và cộng sự \cite{agarwal_2026} mô tả qua phỏng vấn (đáng tin hơn, đổi lại tốn công) hiện ra chủ yếu
ở vế chi phí; vế ``đáng tin hơn'' chỉ thấy ở Codex tầng lớn. Mock và Russo \cite{mock_2026} thấy quy trình tác
tử có thể thêm độ phức tạp không được yêu cầu; ở đây nhánh theo đặc tả có nhiều lỗi tĩnh hơn ở hai ô. Lợi ích
của lập kế hoạch ở mức hàm \cite{jiang_2024} không lặp lại ở mức tính năng khi đề đã chi tiết.

\textbf{Ý nghĩa thực tiễn.} Với người dùng trợ lý lập trình, kết quả gợi ý đầu tư vào đề bài (tên, chữ ký,
kiểu trả về, trường hợp biên) trước khi đầu tư vào quy trình tự động. Nếu dùng quy trình đặc tả, bước có giá
trị nhất có lẽ là để người duyệt bản thiết kế, vì đó là lúc chỗ đề bỏ ngỏ như ở \texttt{f15} lộ ra.

\subsection{Hạn chế và mối đe dọa đến tính hợp lệ}\label{sec:threats}

\textit{Tính hợp lệ cấu trúc.} Để test ẩn chấm được, đề của thí nghiệm chính phải nêu tên và chữ ký, nên nhánh
vibe nhận một yêu cầu đầy đủ hơn vibe coding thật. Thí nghiệm chính đo tác dụng của riêng quy trình khi yêu cầu
đã rõ; E2 được thêm để nhìn vào phần còn lại. Nhánh theo đặc tả là bản rút gọn của Superpowers
(Bảng~\ref{tab:superpowers}): không hỏi đáp với người dùng, không tác tử con, không bước duyệt. Đây chính là
những phần có thể bắt được chỗ đề bỏ ngỏ, và thí nghiệm không đo chúng. Test ẩn cũng kiểm cả chi tiết đề không
chốt, nhưng chúng hạ tỉ lệ đạt ở hai nhánh như nhau nên ít ảnh hưởng hiệu số. Các giới hạn khác của cách đo ở
Phụ lục~\ref{app:validity}.

\textit{Tính hợp lệ nội.} Đề bài do AI soạn sau khi đã biết lời giải. Test ẩn được viết cho bản cài đặt gốc,
vốn làm theo đặc tả, nên có thể thiên vị nhánh này. Tác tử có mạng mở; \texttt{upstream\_urls} chỉ là cận dưới
của việc đọc mã gốc. Trợ lý nghiên cứu cùng mô hình với nhóm \texttt{claude-opus} và đã soạn nháp đề bài mà nhóm
này nhận. Việc làm mù khi gán nhãn chỉ là một phần.

\textit{Tính hợp lệ kết luận.} Tầng vừa có \NumMediumTasks{} và tầng lớn có \NumFeatureTasks{} nhiệm vụ, nên
KTC ở hai tầng này chỉ để tham khảo. Mô phỏng cho độ phủ thật của KTC bootstrap cụm là \NumResCoverageHits{} trên
\NumResCoverageSets{} bộ dữ liệu (\NumResCoveragePct{}\%), thấp hơn mức danh nghĩa, nên KTC có thể hẹp hơn thực
tế; kết cục của H1 và H5 vốn đã là chưa kết luận được nên không đổi. Các phép thử phụ không được hiệu chỉnh so
sánh nhiều lần.

\textit{Tính hợp lệ ngoại.} Kết quả đến từ một dự án, một ngôn ngữ (Python), \NumTasksMain{} nhiệm vụ backend, và
tác giả dự án cũng là thành viên nhóm nghiên cứu. Hai công cụ dùng mô hình khác hãng và khác hạng.

\textit{Quyết định sau khi mở kết quả.} Nhóm \texttt{agy-gemini} bị loại sau khi kết quả chính đã mở, và E1, E2
được thêm sau đó. Cả ba việc được ghi trong nhật ký sai lệch (Phụ lục~\ref{app:deviations}); chúng không đổi
phép thử xác nhận nhưng làm phần thăm dò dễ bị chọn theo kết quả hơn. Trong E2, người chấm máy cùng họ Claude
với ba nhóm và mẫu chấm tay chỉ có sáu lượt Codex, quá ít để kiểm độ nghiêng này.
